% Part: normal-modal-logic
% Chapter: tableaux
% Section: countermodels

\documentclass[../../../include/open-logic-section]{subfiles}

\begin{document}

\olfileid[fr]{nml}{tab}{cou}
      
\olsection{Contre-modèles obtenus à partir de \usetoken{p}{tableau}}

\begin{explain}
  La démonstration du théorème de complétude établit que
  $\Entails !A$ entraîne $\Proves !A$, mais fournit aussi une méthode
  pour construire un contre-modèle de~$!A$ lorsque $\Entails/ !A$.
  Dans~$\Log{K}$, cette méthode constitue une \emph{procédure de
    décision}. Supposons en effet $\Entails/ !A$. La preuve de la
  \olref[cpl]{prop:complete-tableau} indique comment construire un
  !!{tableau} complet, et cette construction se termine toujours.
  Les règles propositionnelles de~$\Log{K}$ n'ajoutent que des
  !!{formula}s préfixées de complexité moindre : pour une formule
  signée $\sFmla{S}{!A}[\sigma]$ donnée, il suffit donc d'appliquer
  chaque règle propositionnelle une seule fois sur une branche. Les
  nouveaux préfixes ne proviennent que de
  \iftag{prvBox}{$\TRule{\False}{\Box}$}{}\iftag{notprvBox,notprvDiamond}{}{
    et }\iftag{prvDiamond}{$\TRule{\True}{\Diamond}$}{}
  \iftag{notprvBox,notprvDiamond}{, règle qu'il suffit}{, règles qu'il suffit}
  d'appliquer une seule fois par formule signée ; chaque application
  crée un seul nouveau préfixe.
  \iftag{prvBox}{$\TRule{\True}{\Box}$}{}\iftag{notprvBox,notprvDiamond}{}{
    et }\iftag{prvDiamond}{$\TRule{\False}{\Diamond}$}{}
  \iftag{notprvBox,notprvDiamond}{peut}{peuvent} devoir s'appliquer
  plusieurs fois, mais une seule fois par préfixe utilisé. Comme il
  n'y a qu'un nombre fini de nouveaux préfixes, la construction
  aboutit en un nombre fini d'étapes à un arbre fermé ou à une
  branche ouverte complète.

  Dès qu'une branche ouverte complète est obtenue, la preuve du
  \olref[cpl]{thm:tableau-completeness} donne un modèle explicite qui
  satisfait l'ensemble initial des !!{formula}s préfixées. Pour décider une conséquence à partir d’un ensemble fini~$\Gamma$, si
  $\Gamma \Entails !A$, il existe un !!{tableau} fermé et
  $\Gamma \Proves !A$. Si la recherche d'un !!{tableau} fermé aboutit plutôt
  à un !!{tableau} complet et ouvert, nous savons que
  $\Gamma \Entails/ !A$ et pouvons construire un contre-modèle.
\end{explain}

\iftag{prvBox}{
\begin{ex}
  Nous savons que $\Proves/ \Box(p \lor q) \lif (\Box p \lor \Box
  q)$. Commençons la construction d'un tableau :
  \begin{oltableau}
    [\pFmla{\False}{\Box(p \lor q) \lif (\Box p \lor \Box q)}{1},
      just = \TAss, checked
      [\pFmla{\True}{\Box(p \lor q)}{1},
        just = {\TRule{\False}{\lif}[1]}, 
        [\pFmla{\False}{\Box p \lor \Box q}{1},
          just = {\TRule{\False}{\lif}[1]}, checked
          [\pFmla{\False}{\Box p}{1},
            just = {\TRule{\False}{\lor}[3]}, checked
            [\pFmla{\False}{\Box q}{1},
              just = {\TRule{\False}{\lor}[3]}, checked
              [\pFmla{\False}{p}{1.1}, 
                just = {\TRule{\False}{\Box}[4]}, checked
                [\pFmla{\False}{q}{1.2}, 
                  just = {\TRule{\False}{\Box}[5]}, checked
                ]
              ]
            ]
          ]
        ]
      ]
    ]
  \end{oltableau}
  Le !!{tableau} n'est pas encore terminé. Examinons la seule ligne
  non cochée : la !!{formula} préfixée
  $\sFmla{\True}{\Box(p \lor q)}[1]$ à la ligne~$2$. La construction
  systématique impose d'appliquer $\TRule{\True}{\Box}$ à chaque
  préfixe successeur déjà utilisé sur cette branche, donc à $1.1$
  et à $1.2$ :
  \begin{oltableau}
    [\pFmla{\False}{\Box(p \lor q) \lif (\Box p \lor \Box q)}{1},
      just = \TAss, checked
      [\pFmla{\True}{\Box(p \lor q)}{1},
        just = {\TRule{\False}{\lif}[1]}, 
        [\pFmla{\False}{\Box p \lor \Box q}{1},
          just = {\TRule{\False}{\lif}[1]}, checked
          [\pFmla{\False}{\Box p}{1},
            just = {\TRule{\False}{\lor}[3]}, checked
            [\pFmla{\False}{\Box q}{1},
              just = {\TRule{\False}{\lor}[3]}, checked
              [\pFmla{\False}{p}{1.1}, 
                just = {\TRule{\False}{\Box}[4]}, checked
                [\pFmla{\False}{q}{1.2}, 
                  just = {\TRule{\False}{\Box}[5]}, checked
                  [\pFmla{\True}{p \lor q}{1.1}, 
                    just = {\TRule{\True}{\Box}[2]}
                    [\pFmla{\True}{p \lor q}{1.2}, 
                      just = {\TRule{\True}{\Box}[2]}
                    ]
                  ]
                ]
              ]
            ]
          ]
        ]
      ]
    ]
  \end{oltableau}
  Les lignes 2, 8 et 9 ne sont pas encore cochées. Aucun nouveau
  préfixe n'ayant été ajouté, appliquons $\TRule{\True}{\lor}$ aux
  lignes~8 et~9 sur toutes les branches qui restent ouvertes :
  \begin{oltableau}
    [\pFmla{\False}{\Box(p \lor q) \lif (\Box p \lor \Box q)}{1},
      just = \TAss, checked
      [\pFmla{\True}{\Box(p \lor q)}{1},
        just = {\TRule{\False}{\lif}[1]}, checked
        [\pFmla{\False}{\Box p \lor \Box q}{1},
          just = {\TRule{\False}{\lif}[1]}, checked
          [\pFmla{\False}{\Box p}{1},
            just = {\TRule{\False}{\lor}[3]}, checked
            [\pFmla{\False}{\Box q}{1},
              just = {\TRule{\False}{\lor}[3]}, checked
              [\pFmla{\False}{p}{1.1}, 
                just = {\TRule{\False}{\Box}[4]}, checked
                [\pFmla{\False}{q}{1.2}, 
                  just = {\TRule{\False}{\Box}[5]}, checked
                  [\pFmla{\True}{p \lor q}{1.1}, 
                    just = {\TRule{\True}{\Box}[2]}, checked
                    [\pFmla{\True}{p \lor q}{1.2}, 
                      just = {\TRule{\True}{\Box}[2]}, checked
                      [\pFmla{\True}{p}{1.1},
                        just = {\TRule{\True}{\lor}[8]}, checked, close
                      ]
                      [\pFmla{\True}{q}{1.1},
                        just = {\TRule{\True}{\lor}[8]}, checked
                        [\pFmla{\True}{p}{1.2},
                          just = {\TRule{\True}{\lor}[9]}, checked]
                        [\pFmla{\True}{q}{1.2},
                          just = {\TRule{\True}{\lor}[9]}, checked, close]
                      ]
                    ]
                  ]
                ]
              ]
            ]
          ]
        ]
      ]
    ]
  \end{oltableau}
  Il reste une branche ouverte, qui est complète. Elle définit un
  modèle dont les mondes sont $W = \{1, 1.1, 1.2\}$, les seuls
  préfixes de cette branche. Sa relation d'accessibilité est
  $R = \{\tuple{1, 1.1}, \tuple{1, 1.2}\}$. Sa valuation est $V(p) =
  \{1.2\}$, car la ligne~11 contient $\sFmla{\True}{p}[1.2]$, et
  $V(q) = \{1.1\}$, car la ligne~10 contient
  $\sFmla{\True}{q}[1.1]$. Le modèle est représenté dans la
  \olref{fig:counter-Box} ; on peut vérifier qu'il réfute
  $\Box(p \lor q) \lif (\Box p \lor \Box q)$.
  \begin{figure}
  \begin{center}
    \begin{tikzpicture}[modal]
      \node[world] (w1) [label={[align=right]right:\mFalse{p}\\ \mFalse{q}}]{$1$}; 
      \node[world] (w2) [label={[align=right]right:\mFalse{p}\\ \mTrue{q}},
        above left=of w1]{$1.1$}; 
      \node[world] (w3) [label={[align=right]right:\mTrue{p}\\ \mFalse{q}},
        above right=of w1] {$1.2$};
      \draw[->] (w1) to (w2);
      \draw[->] (w1) to (w3);
    \end{tikzpicture}
  \end{center}
  \caption{Un contre-modèle de $\Box(p \lor q) \lif (\Box p \lor \Box
    q)$.}
  \ollabel{fig:counter-Box}
\end{figure}
\end{ex}          
}
{
\begin{ex}
  Nous savons que $\Proves/ (\Diamond p \land \Diamond q) \lif \Diamond(p
  \land q)$. Commençons la construction d'un tableau :
  \begin{oltableau}
    [\pFmla{\False}{(\Diamond p \land \Diamond q) \lif \Diamond(p \land q)}{1},
      just = \TAss, checked
      [\pFmla{\True}{\Diamond p \land \Diamond q}{1},
        just = {\TRule{\False}{\lif}[1]}, checked
        [\pFmla{\False}{\Diamond(p \land q)}{1},
          just = {\TRule{\False}{\lif}[1]}
          [\pFmla{\True}{\Diamond p}{1},
            just = {\TRule{\True}{\land}[2]}, checked
            [\pFmla{\True}{\Diamond q}{1},
              just = {\TRule{\True}{\land}[2]}, checked
              [\pFmla{\True}{p}{1.1}, 
                just = {\TRule{\True}{\Diamond}[4]}, checked
                [\pFmla{\True}{q}{1.2}, 
                  just = {\TRule{\True}{\Diamond}[5]}, checked
                ]
              ]
            ]
          ]
        ]
      ]
    ]
  \end{oltableau}
  Le !!{tableau} n'est pas encore terminé. Examinons la seule ligne
  non cochée : la !!{formula} préfixée
  $\sFmla{\False}{\Diamond(p \land q)}[1]$ à la ligne~$3$. La
  construction systématique impose d'appliquer
  $\TRule{\False}{\Diamond}$ à chaque préfixe successeur déjà utilisé
  sur cette branche, donc à $1.1$ et à $1.2$ :
  \begin{oltableau}
    [\pFmla{\False}{(\Diamond p \land \Diamond q) \lif \Diamond(p \land q)}{1},
      just = \TAss, checked
      [\pFmla{\True}{\Diamond p \land \Diamond q}{1},
        just = {\TRule{\False}{\lif}[1]}, checked
        [\pFmla{\False}{\Diamond (p \land q)}{1},
          just = {\TRule{\False}{\lif}[1]}
          [\pFmla{\True}{\Diamond p}{1},
            just = {\TRule{\True}{\land}[2]}, checked
            [\pFmla{\True}{\Diamond q}{1},
              just = {\TRule{\True}{\land}[2]}, checked
              [\pFmla{\True}{p}{1.1}, 
                just = {\TRule{\True}{\Diamond}[4]}, checked
                [\pFmla{\True}{q}{1.2}, 
                  just = {\TRule{\True}{\Diamond}[5]}, checked
                  [\pFmla{\False}{p \land q}{1.1}, 
                    just = {\TRule{\False}{\Diamond}[3]}
                    [\pFmla{\False}{p \land q}{1.2}, 
                      just = {\TRule{\False}{\Diamond}[3]}
                    ]
                  ]
                ]
              ]
            ]
          ]
        ]
      ]
    ]
  \end{oltableau}
  Les lignes 3, 8 et 9 ne sont pas encore cochées. Aucun nouveau
  préfixe n'ayant été ajouté, appliquons $\TRule{\False}{\land}$ aux
  lignes~8 et~9 sur toutes les branches qui restent ouvertes :
  \begin{oltableau}
    [\pFmla{\False}{(\Diamond p \land \Diamond q) \lif \Diamond(p \land q)}{1},
      just = \TAss, checked
      [\pFmla{\True}{\Diamond p \land \Diamond q}{1},
        just = {\TRule{\False}{\lif}[1]}, checked
        [\pFmla{\False}{\Diamond(p \land q)}{1},
          just = {\TRule{\False}{\lif}[1]}, checked
          [\pFmla{\True}{\Diamond p}{1},
            just = {\TRule{\True}{\land}[2]}, checked
            [\pFmla{\True}{\Diamond q}{1},
              just = {\TRule{\True}{\land}[2]}, checked
              [\pFmla{\True}{p}{1.1}, 
                just = {\TRule{\True}{\Diamond}[4]}, checked
                [\pFmla{\True}{q}{1.2}, 
                  just = {\TRule{\True}{\Diamond}[5]}, checked
                  [\pFmla{\False}{p \land q}{1.1}, 
                    just = {\TRule{\False}{\Diamond}[3]}, checked
                    [\pFmla{\False}{p \land q}{1.2}, 
                      just = {\TRule{\False}{\Diamond}[3]}, checked
                      [\pFmla{\False}{p}{1.1},
                        just = {\TRule{\False}{\land}[8]}, checked, close
                      ]
                      [\pFmla{\False}{q}{1.1},
                        just = {\TRule{\False}{\land}[8]}, checked
                        [\pFmla{\False}{p}{1.2},
                          just = {\TRule{\False}{\land}[9]}, checked]
                        [\pFmla{\False}{q}{1.2},
                          just = {\TRule{\False}{\land}[9]}, checked, close]
                      ]
                    ]
                  ]
                ]
              ]
            ]
          ]
        ]
      ]
    ]
  \end{oltableau}
  Il reste une branche ouverte, qui est complète. Elle définit un
  modèle dont les mondes sont $W = \{1, 1.1, 1.2\}$, les seuls
  préfixes de cette branche. Sa relation d'accessibilité est
  $R = \{\tuple{1, 1.1}, \tuple{1, 1.2}\}$. Sa valuation est $V(p) =
  \{1.1\}$, car la ligne~6 contient $\sFmla{\True}{p}[1.1]$, et
  $V(q) = \{1.2\}$, car la ligne~7 contient
  $\sFmla{\True}{q}[1.2]$. Le modèle est représenté dans la
  \olref{fig:counter-Diamond} ; on peut vérifier qu'il réfute
  $(\Diamond p \land \Diamond q) \lif \Diamond (p \land q)$.
  \begin{figure}
  \begin{center}
    \begin{tikzpicture}[modal]
      \node[world] (w1) [label={[align=right]right:\mFalse{p}\\ \mFalse{q}}]{$1$}; 
      \node[world] (w2) [label={[align=right]right:\mTrue{p}\\ \mFalse{q}},
        above left=of w1]{$1.1$}; 
      \node[world] (w3) [label={[align=right]right:\mFalse{p}\\ \mTrue{q}},
        above right=of w1] {$1.2$};
      \draw[->] (w1) to (w2);
      \draw[->] (w1) to (w3);
    \end{tikzpicture}
  \end{center}
  \caption{Un contre-modèle de $(\Diamond p \land \Diamond q) \lif
    \Diamond (p \land q)$.}
  \ollabel{fig:counter-Diamond}
\end{figure}
\end{ex}          
}
\begin{explain}
Dans les deux exemples, les mondes du contre-modèle sont
les trois préfixes de la branche ouverte complète. Pour
déterminer la valuation d'une variable, on relève les
préfixes auxquels cette variable figure avec le signe
$\True$ sur la branche. Les formules signées au préfixe~$1$
donnent alors la valeur des hypothèses et de la conclusion
au monde initial.
\end{explain}
\end{document}
